% GENERATED FILE. Edit canonical lessons, then run npm run generate:books.
% canonical-source-hash: fnv1a64:c328040720b218fb
% canonical-lessons: KA-C289-hegadaru, KA-C289-andare, KA-C289-nepa, KA-C289-adhara, KA-C289-saksi, KA-R289-first-pass-words-from-chapters-281-284, KA-R289-second-pass-words-from-chapters-285-289

\chapter{Somehow, That is to say, Excuse, Basis, Evidence}
\label{ch:ka-somehow-that-is-to-say-excuse-basis-evidence}
\hlchaptermodality{289}

\begin{chapteropening}
\textbf{By the end of this chapter:} \emph{I can say somehow, that is to say, an excuse, a basis and evidence.}

Complete the last lesson of chapter 289: I can say somehow, that is to say, an excuse, a basis and evidence.
\end{chapteropening}

\section[\texorpdfstring{\kn{ಹೇಗಾದರೂ} (hēgādarū)}{ಹೇಗಾದರೂ (hēgādarū)}]{\kn{ಹೇಗಾದರೂ} (hēgādarū) — somehow, one way or another}

\label{lesson:KA-C289-hegadaru}



\begin{quote}
Before the new one: say the Kannada for without, then the Kannada for anyway.
\end{quote}

\subsection*{You'll want to know: \kn{ಹೇಗಾದರೂ}}
\textbf{\kn{ಹೇಗಾದರೂ}} — \emph{hēgādarū} — \textquotedblleft{}somehow, one way or another\textquotedblright{}.

\begin{grammarlens}[title={What you've built}]
One more word for this chapter.
\end{grammarlens}

\subsection*{Your turn}
\begin{itemize}
\raggedright
  \item \emph{Say it:} \emph{hēgādarū}
  \item \emph{Say it:} \emph{hēgādarū}, once more
  \item \emph{Say it:} say \emph{hēgādarū}
  \item \emph{Recall:} say the Kannada for without, then the Kannada for anyway, then say \emph{hēgādarū} again
\end{itemize}

\begin{tcolorbox}[breakable,colback=teal!4,colframe=teal!35!black,title={Before you move on}]
What does \kn{ಹೇಗಾದರೂ} mean? (\textquotedblleft{}Somehow, one way or another\textquotedblright{}.) Say it once more.
\end{tcolorbox}
\section[\texorpdfstring{\kn{ಅಂದರೆ} (andare)}{ಅಂದರೆ (andare)}]{\kn{ಅಂದರೆ} (andare) — that is to say, meaning}

\label{lesson:KA-C289-andare}



\begin{quote}
Before the new one: say the Kannada for anyway, then the Kannada for somehow.
\end{quote}

\subsection*{You'll want to know: \kn{ಅಂದರೆ}}
\textbf{\kn{ಅಂದರೆ}} — \emph{andare} — \textquotedblleft{}that is to say, meaning\textquotedblright{}.

\begin{grammarlens}[title={What you've built}]
One more word for this chapter.
\end{grammarlens}

\subsection*{Your turn}
\begin{itemize}
\raggedright
  \item \emph{Say it:} \emph{andare}
  \item \emph{Say it:} \emph{andare}, once more
  \item \emph{Say it:} say \emph{andare}
  \item \emph{Recall:} say the Kannada for anyway, then the Kannada for somehow, then say \emph{andare} again
\end{itemize}

\begin{tcolorbox}[breakable,colback=teal!4,colframe=teal!35!black,title={Before you move on}]
What does \kn{ಅಂದರೆ} mean? (\textquotedblleft{}That is to say, meaning\textquotedblright{}.) Say it once more.
\end{tcolorbox}
\section[\texorpdfstring{\kn{ನೆಪ} (nepa)}{ನೆಪ (nepa)}]{\kn{ನೆಪ} (nepa) — an excuse, a pretext}

\label{lesson:KA-C289-nepa}



\begin{quote}
Before the new one: say the Kannada for somehow, then the Kannada for that is to say.
\end{quote}

\subsection*{You'll want to know: \kn{ನೆಪ}}
\textbf{\kn{ನೆಪ}} — \emph{nepa} — \textquotedblleft{}an excuse, a pretext\textquotedblright{}.

\begin{grammarlens}[title={What you've built}]
One more word for this chapter.
\end{grammarlens}

\subsection*{Your turn}
\begin{itemize}
\raggedright
  \item \emph{Say it:} \emph{nepa}
  \item \emph{Say it:} \emph{nepa}, once more
  \item \emph{Say it:} say \emph{nepa}
  \item \emph{Recall:} say the Kannada for somehow, then the Kannada for that is to say, then say \emph{nepa} again
\end{itemize}

\begin{tcolorbox}[breakable,colback=teal!4,colframe=teal!35!black,title={Before you move on}]
What does \kn{ನೆಪ} mean? (\textquotedblleft{}An excuse, a pretext\textquotedblright{}.) Say it once more.
\end{tcolorbox}
\section[\texorpdfstring{\kn{ಆಧಾರ} (ādhāra)}{ಆಧಾರ (ādhāra)}]{\kn{ಆಧಾರ} (ādhāra) — a basis, grounds; support}

\label{lesson:KA-C289-adhara}



\begin{quote}
Before the new one: say the Kannada for that is to say, then the Kannada for an excuse.
\end{quote}

\subsection*{You'll want to know: \kn{ಆಧಾರ}}
\textbf{\kn{ಆಧಾರ}} — \emph{ādhāra} — \textquotedblleft{}a basis, grounds; support\textquotedblright{}.

\begin{grammarlens}[title={What you've built}]
One more word for this chapter.
\end{grammarlens}

\subsection*{Your turn}
\begin{itemize}
\raggedright
  \item \emph{Say it:} \emph{ādhāra}
  \item \emph{Say it:} \emph{ādhāra}, once more
  \item \emph{Say it:} say \emph{ādhāra}
  \item \emph{Recall:} say the Kannada for that is to say, then the Kannada for an excuse, then say \emph{ādhāra} again
\end{itemize}

\begin{tcolorbox}[breakable,colback=teal!4,colframe=teal!35!black,title={Before you move on}]
What does \kn{ಆಧಾರ} mean? (\textquotedblleft{}A basis, grounds; support\textquotedblright{}.) Say it once more.
\end{tcolorbox}
\section[\texorpdfstring{\kn{ಸಾಕ್ಷಿ} (sākṣi)}{ಸಾಕ್ಷಿ (sākṣi)}]{\kn{ಸಾಕ್ಷಿ} (sākṣi) — evidence; a witness}

\label{lesson:KA-C289-saksi}



\begin{quote}
Before the new one: say the Kannada for an excuse, then the Kannada for a basis.
\end{quote}

\subsection*{You'll want to know: \kn{ಸಾಕ್ಷಿ}}
\textbf{\kn{ಸಾಕ್ಷಿ}} — \emph{sākṣi} — \textquotedblleft{}evidence; a witness\textquotedblright{}.

\begin{grammarlens}[title={What you've built}]
One more word for this chapter.
\end{grammarlens}

\subsection*{Your turn}
\begin{itemize}
\raggedright
  \item \emph{Say it:} \emph{sākṣi}
  \item \emph{Say it:} \emph{sākṣi}, once more
  \item \emph{Say it:} say \emph{sākṣi}
  \item \emph{Recall:} say the Kannada for an excuse, then the Kannada for a basis, then say \emph{sākṣi} again
\end{itemize}

\begin{tcolorbox}[breakable,colback=teal!4,colframe=teal!35!black,title={Before you move on}]
What does \kn{ಸಾಕ್ಷಿ} mean? (\textquotedblleft{}Evidence; a witness\textquotedblright{}.) Say it once more.
\end{tcolorbox}
\section[(dialogue)]{Words from chapters 281-284 — a first pass}

\label{lesson:KA-R289-first-pass-words-from-chapters-281-284}



\begin{quote}
Say the words for a basis and evidence.
\end{quote}

\subsection*{Your turn}
Say these aloud:

\begin{itemize}
\raggedright
  \item to vanish, to pack, to take off, to fail, to get stuck
  \item an accident, an ambulance, stress, a nurse, surprise
  \item the pulse, a cramp, a fracture, an operation, pregnant
  \item full moon, sunset, yesterday, time, sometimes
\end{itemize}

\begin{tcolorbox}[breakable,colback=teal!4,colframe=teal!35!black,title={Before you move on}]
To vanish? (\textbf{\kn{ಮಾಯವಾಗು}}, \emph{māyavāgu}.) To pack? (\textbf{\kn{ಪ್ಯಾಕ್} \kn{ಮಾಡು}}, \emph{pyāk māḍu}.) To take off? (\textbf{\kn{ಕಳಚು}}, \emph{kaḷacu}.) To fail? (\textbf{\kn{ಫೇಲಾಗು}}, \emph{phēlāgu}.) To get stuck? (\textbf{\kn{ಸಿಕ್ಕಿಹಾಕಿಕೊಳ್ಳು}}, \emph{sikkihākikoḷḷu}.) An accident? (\textbf{\kn{ಅಪಘಾತ}}, \emph{apaghāta}.) An ambulance? (\textbf{\kn{ಆಂಬ್ಯುಲೆನ್ಸ್}}, \emph{āmbyulens}.) Stress? (\textbf{\kn{ಒತ್ತಡ}}, \emph{ottaḍa}.) A nurse? (\textbf{\kn{ದಾದಿ}}, \emph{dādi}.) Surprise? (\textbf{\kn{ಆಶ್ಚರ್ಯ}}, \emph{āścarya}.) The pulse? (\textbf{\kn{ನಾಡಿ}}, \emph{nāḍi}.) A cramp? (\textbf{\kn{ಸೆಳೆತ}}, \emph{seḷeta}.) A fracture? (\textbf{\kn{ಮುರಿತ}}, \emph{murita}.) An operation? (\textbf{\kn{ಆಪರೇಷನ್}}, \emph{āpareṣan}.) Pregnant? (\textbf{\kn{ಗರ್ಭಿಣಿ}}, \emph{garbhiṇi}.) Full moon? (\textbf{\kn{ಹುಣ್ಣಿಮೆ}}, \emph{huṇṇime}.) Sunset? (\textbf{\kn{ಸೂರ್ಯಾಸ್ತ}}, \emph{sūryāsta}.) Yesterday? (\textbf{\kn{ನಿನ್ನೆ}}, \emph{ninne}.) Time? (\textbf{\kn{ಸಮಯ}}, \emph{samaya}.) Sometimes? (\textbf{\kn{ಕೆಲವೊಮ್ಮೆ}}, \emph{kelavomme}.)
\end{tcolorbox}
\section[(dialogue)]{Words from chapters 285-289 — a second pass}

\label{lesson:KA-R289-second-pass-words-from-chapters-285-289}



\begin{quote}
Say the words for three quarters and the next day.
\end{quote}

\subsection*{Your turn}
Say these aloud:

\begin{itemize}
\raggedright
  \item three quarters, the next day, childhood, the new moon day, an interval
  \item a building, a floor, an entrance, a platform, abroad
  \item a riverbank, a drought, an earthquake, the world, a wall
  \item moreover, suddenly, except, without, anyway
  \item somehow, that is to say, an excuse, a basis, evidence
\end{itemize}

\begin{tcolorbox}[breakable,colback=teal!4,colframe=teal!35!black,title={Before you move on}]
Three quarters? (\textbf{\kn{ಮುಕ್ಕಾಲು}}, \emph{mukkālu}.) The next day? (\textbf{\kn{ಮರುದಿನ}}, \emph{marudina}.) Childhood? (\textbf{\kn{ಬಾಲ್ಯ}}, \emph{bālya}.) The new moon day? (\textbf{\kn{ಅಮಾವಾಸ್ಯೆ}}, \emph{amāvāsye}.) An interval? (\textbf{\kn{ಮಧ್ಯಂತರ}}, \emph{madhyantara}.) A building? (\textbf{\kn{ಕಟ್ಟಡ}}, \emph{kaṭṭaḍa}.) A floor? (\textbf{\kn{ಮಹಡಿ}}, \emph{mahaḍi}.) An entrance? (\textbf{\kn{ಪ್ರವೇಶ}}, \emph{pravēśa}.) A platform? (\textbf{\kn{ವೇದಿಕೆ}}, \emph{vēdike}.) Abroad? (\textbf{\kn{ವಿದೇಶ}}, \emph{vidēśa}.) A riverbank? (\textbf{\kn{ದಡ}}, \emph{daḍa}.) A drought? (\textbf{\kn{ಬರ}}, \emph{bara}.) An earthquake? (\textbf{\kn{ಭೂಕಂಪ}}, \emph{bhūkampa}.) The world? (\textbf{\kn{ಪ್ರಪಂಚ}}, \emph{prapañca}.) A wall? (\textbf{\kn{ಗೋಡೆ}}, \emph{gōḍe}.) Moreover? (\textbf{\kn{ಅಲ್ಲದೆ}}, \emph{allade}.) Suddenly? (\textbf{\kn{ಇದ್ದಕ್ಕಿದ್ದಂತೆ}}, \emph{iddakkiddante}.) Except? (\textbf{\kn{ಹೊರತು}}, \emph{horatu}.) Without? (\textbf{\kn{ಇಲ್ಲದೆ}}, \emph{illade}.) Anyway? (\textbf{\kn{ಹೇಗೂ}}, \emph{hēgū}.) Somehow? (\textbf{\kn{ಹೇಗಾದರೂ}}, \emph{hēgādarū}.) That is to say? (\textbf{\kn{ಅಂದರೆ}}, \emph{andare}.) An excuse? (\textbf{\kn{ನೆಪ}}, \emph{nepa}.) A basis? (\textbf{\kn{ಆಧಾರ}}, \emph{ādhāra}.) Evidence? (\textbf{\kn{ಸಾಕ್ಷಿ}}, \emph{sākṣi}.)
\end{tcolorbox}
